\chapter{Disaggregated Serving}\label{sec:L6}

This chapter extends the tools of \cref{sec:L1,sec:L2,sec:L3,sec:L4} and the colocated
scheduler of \cref{sec:L5} to a disaggregated deployment. There, prefill and decode run
on separate instances with separate memory pools, and a link carries the KV cache
between them. The chapter derives what a scheduler for it needs: what a miss costs at
each place where it can happen, what each instance should keep in memory and for how
long, how many sessions to admit, and when misses feed themselves.

\section{The disaggregated replica}\label{L6:sec:model}

\Cref{fig:deployment} shows the model.

\begin{figure}[t]
\centering
\begin{tikzpicture}[x=0.8cm, y=0.8cm, >=Stealth, font=\small, line width=0.6pt,
  station/.style={circle, draw, thick, minimum size=10mm, align=center, inner sep=1pt},
  note/.style={font=\footnotesize, align=center, text=black!75}]
  % ---------- prefill instance ----------
  \foreach \i in {0,1}
    \draw[fill=orange!10] (0.4+0.5*\i,-0.3) rectangle ++(0.45,0.6);
  \draw (0.35,-0.3) -- (1.4,-0.3) (0.35,0.3) -- (1.4,0.3);
  \node[draw, thick, rounded corners=1pt, fill=orange!10, minimum width=12mm, minimum height=9mm] (SP) at (2.5,0) {};
  \foreach \x in {-0.36,0,0.36}
    \foreach \y in {-0.17,0.17}
      \draw[orange!60!black, fill=white] ($(SP)+(\x-0.13,\y-0.11)$) rectangle ++(0.26,0.22);
  \node[station, fill=blue!8] (P) at (4.4,0) {FIFO};
  \draw[->] (SP.east) -- (P.west);
  \draw[fill=gray!15, draw=gray!60] (0.2,-1.4) rectangle (5.3,-0.85);
  \node[font=\scriptsize, text=black!75] at (2.75,-1.12) {prefix cache of paused sessions};
  \node[note] at (2.5,0.85) {$m_P$ slots};
  \node[note] at (4.4,0.9) {prefill};
  \node[draw, dashed, rounded corners=3mm, inner sep=2.5mm, fit={(0.2,-1.4) (5.3,1.2)}] (MP) {};
  \node[note, anchor=south] at (MP.north) {\textbf{Prefill instance} (pool $M_P$)};
  % ---------- link ----------
  \node[station, fill=green!8] (L) at (7.0,0) {PS};
  \node[note, anchor=north] at (7.0,-0.65) {KV link\\bandwidth $B$};
  % ---------- decode instance ----------
  \foreach \i in {0,1}
    \draw[fill=orange!10] (8.9+0.5*\i,-0.3) rectangle ++(0.45,0.6);
  \draw (8.85,-0.3) -- (9.9,-0.3) (8.85,0.3) -- (9.9,0.3);
  \node[draw, thick, rounded corners=1pt, fill=orange!10, minimum width=12mm, minimum height=9mm] (SD) at (11.0,0) {};
  \foreach \x in {-0.36,0,0.36}
    \foreach \y in {-0.17,0.17}
      \draw[orange!60!black, fill=white] ($(SD)+(\x-0.13,\y-0.11)$) rectangle ++(0.26,0.22);
  \node[station, fill=red!8] (D) at (12.9,0) {PS};
  \draw[->] (SD.east) -- (D.west);
  \node[note] at (11.0,0.85) {$m_D$ slots};
  \node[note] at (12.9,-0.95) {batch $\varphi(m)$};
  \node[draw, dashed, rounded corners=3mm, inner sep=2.5mm, fit={(8.8,-1.25) (14.0,1.2)}] (MD) {};
  \node[note, anchor=south] at (MD.north) {\textbf{Decode instance} (pool $M_D$)};
  % ---------- tool call ----------
  \node[draw, thick, rounded corners, fill=teal!8, minimum width=20mm, minimum height=9mm] (T) at (7.0,-3.1) {};
  \foreach \k in {-2,-1,0,1,2}
    \draw[teal!60!black] ($(T)+(0.36*\k,0)$) circle (0.13);
  \node[note, anchor=north] at ($(T.south)+(0,-0.05)$) {\textbf{Tool call}: delay station};
  % ---------- flows ----------
  \draw[->] (-3.0,0) -- (0.3,0);
  \node[note, anchor=south] at (-2.1,0.05) {new sessions\\Poisson $\Lambda$};
  \draw[->] (P.east) -- (L.west);
  \draw[->] (L.east) -- (8.8,0);
  \draw[->] (D.east) -- (15.4,0) node[note, right] {ends};
  \draw[->] (D.south) |- node[note, right, pos=0.25] {resume w.p.\ $p$} (T.east);
  \draw[->] (T.west) -| (-0.2,0) -- (0.3,0);
  \node[note, anchor=east] at (-0.3,-2.2) {next turn\\(hit or miss)};
\end{tikzpicture}
\caption{A disaggregated deployment as a queueing model. A turn waits for KV memory
at a prefill instance (it is admitted only when its whole context fits, in arrival
order), is prefilled one request at a time in arrival order, sends its KV cache over a
link whose bandwidth concurrent transfers share, waits for KV memory at a decode
instance, and decodes in a batch that shares the instance's throughput. Sessions loop
through a tool call back to the prefill instance, where the prefix cache decides
whether the next turn is a hit (short prefill work) or a miss (the whole context is
recomputed).}
\label{fig:deployment}
\end{figure}

\begin{definition}[Workload and costs]\label{L6:def:workload}
Sessions arrive as a Poisson process of rate $\Lambda$ (\cref{def:poisson}). Turn $j$
of session $i$ has a cached prefix of $K_{ij}$ tokens, $n_{ij}$ new tokens and
$o_{ij}$ output tokens, with context $T_{ij}=K_{ij}+n_{ij}$ and
\begin{equation}\label{L6:eq:context}
  K_{i1}=0,\qquad K_{i,j+1}=T_{ij},\qquad n_{i,j+1}\ge o_{ij},
\end{equation}
since the next turn's new tokens contain the previous output and the tool's output.
After each turn the session calls a tool for a time with mean $Z$. A replica has
prefill constants $a,b>0$, $\kappa$ KV bytes per token, a link of bandwidth $B$ bytes
per second with overhead $x_0$, a decode capacity $\varphi$, and pools of $M_P$ and
$M_D$ bytes. With the hit indicator $H_{ij}$ (the prefix is still cached at
admission), turn $(i,j)$ has
\begin{equation}\label{L6:eq:work}
  S_{ij}=P(n_{ij},K_{ij})+(1-H_{ij})\,\Delta S(K_{ij}),\qquad
  X_{ij}=x_0+\frac{\kappa T_{ij}}{B},\qquad D_{ij}=o_{ij},\qquad c_{ij}=\kappa T_{ij},
\end{equation}
the prefill work in seconds ($P$ and $\Delta S(K)=aK+bK^2/2$ as in
\cref{L5:def:replica,cor:price-mono}), the link work in seconds, the decode work in
tokens, and the footprint in bytes on either instance. A batch of $m$ decodes takes
$\tau(m)$ seconds per step, so $\varphi(m)=m/\tau(m)$ tokens per second.
\end{definition}

\begin{definition}[Path of a turn]\label{L6:def:path}
\begin{enumerate}[label=(\roman*)]
\item A turn waits in arrival order at the prefill pool until $c=\kappa T$ bytes of
  running turns' allocations are free. Cached bytes are evicted as needed and never
  block admission. At admission the turn uses its own cached prefix if it survived;
  a hit is all or nothing.
\item It receives $S$ seconds of FIFO prefill, which produces the first token, and
  $X$ seconds of link work. It then releases its prefill allocation and leaves its
  $\kappa T$ bytes cached. Ending a session does not erase this prefix.
\item It waits in the same way for $c$ bytes at the decode pool, receives $D$ tokens of
  processor-sharing decode, and releases the allocation.
\item With probability $p<1$ the session calls a tool and submits another turn;
  otherwise it ends. So the number of turns $J$ has $\E[J]=1/(1-p)$.
\end{enumerate}
A FIFO stage serves only its first turn, at rate one; a PS stage with $m$ turns serves
each at rate $\varphi(m)/m$; the tool serves at rate one without waiting. Allocated plus
cached bytes never exceed a pool.
\end{definition}

The executable transfer of \cref{app:serq} reserves the decode pool before the
transfer, which changes how long each pool is held.

\section{Sample paths}\label{L6:sec:paths}

Number the turns in order of arrival at the prefill pool, and let turn $k$ of session
$\sigma(k)$ have the epochs
$a_k\le\alpha_k\le s_k\le f_k\le e_k\le\alpha'_k\le d_k$: arrival, admission, start and
end of prefill, end of transfer, admission at decode, and end of decode. Let
$R_P(t)=\{k:\alpha_k\le t<e_k\}$ be the turns holding prefill memory. Then
(\cref{L6:ex:paths})
\begin{align}
  \alpha_k&=\inf\Bigl\{t\ge\max(a_k,\alpha_{k-1}):
      \textstyle\sum_{l\in R_P(t)}c_l+c_k\le M_P\Bigr\},\label{L6:eq:admitP}\\
  s_k&=\max(\alpha_k,f_{k-1}),\qquad f_k=s_k+S_k,\label{L6:eq:lindley}\\
  a_{k^+}&=d_k+Z_k\quad\text{if the session continues},\label{L6:eq:feedback}
\end{align}
where $k^+$ is the next turn of $\sigma(k)$. \Cref{L6:eq:lindley} is the FIFO recursion
of \cref{sec:L3,sec:L4}, and \eqref{L6:eq:feedback} is the closed loop. If the system
is stable, the turn rate is still fixed by the session rate,
$\lambda=\Lambda\E[J]=\Lambda/(1-p)$. The link and decode batch are PS stations, the
tool a delay station, and each pool a slot queue (\cref{L2:sec:mmm}) with
$m_j=\lfloor M_j/c\rfloor$ slots; a prefill slot is held through the transfer, until
$e_k$.

\paragraph{The cache chooses the prefill work.} The eviction order decides $H_k$, so
with hit rate $h=\E[H]$
\begin{equation}\label{L6:eq:ES}
  \E[S]=\E\bigl[P(n,K)\bigr]+\E\bigl[(1-H)\,\Delta S(K)\bigr],
\end{equation}
and only the second term depends on the policy. \Cref{L5:lem:split} splits the
variance of this work.

\paragraph{Performance.} The first token appears when the prefill ends, so
\begin{equation}\label{L6:eq:ttft}
  \mathrm{TTFT}_k=f_k-a_k
  ={\underbrace{(\alpha_k-a_k)}_{W^M_{P,k}}}
  +{\underbrace{(s_k-\alpha_k)}_{W_{P,k}}}+S_k ,
\end{equation}
the wait for prefill memory, the wait for the prefill server ($W_P$ is its mean), and
the prefill work. Transfer and decode reach the TTFT only through memory, because a
prefill slot is held until $e_k$. A client that sees its first token only after the
handoff also waits for the link and the decode pool. No stage is stable unless
\begin{equation}\label{L6:eq:stable}
  \rho_P=\lambda\,\E[S]<1,\qquad \rho_L=\lambda\,\E[X]<1,\qquad
  \lambda\,\E[D]<\max_{m\le m_D}\varphi(m),
\end{equation}
and by \eqref{L6:eq:ES} the first condition already depends on the eviction policy.

\section{Three prices}\label{L6:sec:prices}

A decision that changes where a paused session's KV lives changes the work of its next
turn at one of three stations.

\paragraph{(a) A miss at the prefill instance.} Under independent Poisson arrivals and i.i.d. prefill work, the
prefill instance is an M/G/1 FIFO queue, so \cref{thm:price} applies as it
stands: turning a fraction $q_i$ of arrivals from hits into misses raises the mean number
at the prefill instance by an amount in
$[\lambda\sum_iq_i\Phi_i,\ \frac{1-\rho_P}{1-\rho_P'}\lambda\sum_iq_i\Phi_i]$, with
\[
  \Phi_i=\Delta S_i
   +\frac{\lambda\bigl((S_i^{\mathrm{miss}})^2-(S_i^{\mathrm{hit}})^2\bigr)}{2(1-\rho_P)}
   +\frac{\lambda W_P\,\Delta S_i}{1-\rho_P}.
\]
The middle term is the head-of-line term: every turn that arrives while the miss is
prefilled waits for all of it.

\paragraph{(b) A miss at the decode instance.} The path of \cref{L6:def:path} sends
the whole context every turn. Suppose now that the decode instance also keeps the KV
of paused sessions. If it still holds the
session's KV from its previous turn, the next turn transfers only its new tokens,
$T_i=n_i$; if not, it transfers the whole context, $T_i=K_i+n_i$. A decode-side miss
therefore adds $\Delta X_i=\kappa K_i/B$ of link work.

\begin{proposition}[The price of a transfer miss]\label{L6:prop:link-price}
Let the link be a processor-sharing station of capacity 1 with load $\rho_L<1$, and
turn a fraction $q_i$ of arrivals into decode-side misses, raising the load to
$\rho_L'=\rho_L+\lambda\sum_iq_i\Delta X_i<1$. Then the mean number of transfers in
progress rises exactly by
\[
  \Delta L_L=\frac{1-\rho_L}{1-\rho_L'}\;\lambda\sum_iq_i\,\Phi^X_i,\qquad
  \Phi^X_i=\frac{\Delta X_i}{(1-\rho_L)^2},
\]
whatever the distribution of the transfer times.
\end{proposition}

\begin{proof}
By insensitivity (\cref{thm:insens}) the number of transfers in progress has the
geometric law of \cref{prop:decode-price} with capacity $C=1$, for every transfer-time
distribution with the given mean. \Cref{prop:decode-price} with $C=1$ and $\Delta D_i$
replaced by $\Delta X_i$ gives the formula.
\end{proof}

\paragraph{(c) Decode demand.} A decision that adds decode demand, such as admitting a
session, pays the load factor of \cref{prop:decode-price} at the decode batch; a miss at
either instance adds none (\cref{L2:cor:decode-hit}).

\paragraph{Comparison.} The same context of $K$ tokens costs the prefill instance a
miss whose work is quadratic in $K$ and whose price has a head-of-line term, set by the
second moment of the prefill work; it costs the link a transfer linear in $K$ whose
price depends on its mean only, because under processor sharing a long transfer does not
block the short ones. This growth comparison does not rank the pools universally: bandwidth,
load, resume probability and holding time also set their memory values.

\section{The decision problem with two pools}\label{L6:sec:decision}

At a decision epoch each instance holds KV for a set of paused sessions. Session $i$
holds $c_i$ bytes, resumes with probability $p_i$ after an expected further $\tau_i$
seconds, and its next turn pays $\Phi_i$ if its prefix is not in the prefill instance's
cache and $\Phi^X_i$ if its KV is not at the decode instance. With indicators
$x^P_i,x^D_i\in\{0,1\}$ of ``keep'' at each instance and byte-second budgets $\mathcal B_P$ and
$\mathcal B_D$, the keep-or-drop problem is
\begin{equation}\label{L6:eq:primal}
  \mathrm{OPT}=\min_{x^P,x^D}\ \sum_i\bigl[(1-x^P_i)\,p_i\Phi_i+(1-x^D_i)\,p_i\Phi^X_i\bigr]
  \quad\text{s.t.}\quad \sum_ic_i\tau_ix^P_i\le \mathcal B_P,\ \ \sum_ic_i\tau_ix^D_i\le \mathcal B_D .
\end{equation}
Relaxing the two constraints with multipliers $\theta_P,\theta_D\ge0$, the
\term{shadow prices} of a byte-second at each instance, gives the Lagrangian
$\mathcal{L}=\sum_i[(1-x^P_i)p_i\Phi_i+(1-x^D_i)p_i\Phi^X_i]
+\theta_P(\sum_ic_i\tau_ix^P_i-\mathcal B_P)+\theta_D(\sum_ic_i\tau_ix^D_i-\mathcal B_D)$.

\begin{proposition}[One price per pool]\label{L6:prop:shadow}
Let $u^P_i=p_i\Phi_i/(c_i\tau_i)$ and $u^X_i=p_i\Phi^X_i/(c_i\tau_i)$.
\begin{enumerate}[label=(\roman*)]
\item For all $\theta_P,\theta_D\ge0$, $\mathcal{L}$ is minimised by keeping session $i$ at
  the prefill instance iff $u^P_i>\theta_P$ and at the decode instance iff
  $u^X_i>\theta_D$, and the minimum $g(\theta_P,\theta_D)$ is at most $\mathrm{OPT}$.
\item If the multipliers are such that these choices use both budgets exactly, the
  choices are optimal for \eqref{L6:eq:primal}.
\end{enumerate}
\end{proposition}

\begin{proof}
(i) $\mathcal{L}=\sum_ip_i(\Phi_i+\Phi^X_i)-\theta_P\mathcal B_P-\theta_D\mathcal B_D
+\sum_ix^P_i(\theta_Pc_i\tau_i-p_i\Phi_i)+\sum_ix^D_i(\theta_Dc_i\tau_i-p_i\Phi^X_i)$ is
separable: each $x^P_i$ is best set to 1 exactly when $\theta_Pc_i\tau_i<p_i\Phi_i$, that
is $u^P_i>\theta_P$, and likewise each $x^D_i$. For a feasible $(x^P,x^D)$ both brackets
multiplying $\theta_P$ and $\theta_D$ are nonpositive, so its objective is at least its
Lagrangian, which is at least $g$; minimising over feasible choices gives
$\mathrm{OPT}\ge g$ (weak duality). (ii) The choices of (i) are then feasible and their
constraint terms vanish, so their objective equals $g\le\mathrm{OPT}$.
\end{proof}

This is the threshold rule of \cref{thm:threshold}, once per pool: each instance keeps
what is worth more per byte-second than its own memory.

\begin{corollary}[Free transfer]\label{L6:cor:free-transfer}
As the link bandwidth $B\to\infty$, $\Phi^X_i\to0$ for every session, so for any $\theta_D>0$ the decode
instance keeps nothing for paused sessions, and all of $M_D$ serves running decodes.
\end{corollary}

\begin{proof}
$\Delta X_i=\kappa K_i/B\to0$, and $\rho_L\to\lambda x_0$ stays below 1 when $\lambda x_0<1$,
so $\Phi^X_i=\Delta X_i/(1-\rho_L)^2\to0$ and $u^X_i\to0<\theta_D$.
\end{proof}

A decode instance's cache is worth exactly the transfers it saves; the prefill
instance's cache is worth the recomputation it saves and the queueing that recomputation
would cause.

\section{How long to keep}\label{L6:sec:hold}
The holding-time proof in \cref{L5:prop:hold} applies separately to each pool.
Use $(\Phi_i,\theta_P)$ at prefill and $(\Phi^X_i,\theta_D)$ at decode.
The resume hazard, the exponential-gap holding time, and the unknown-gap
ski-rental comparison are unchanged. They hold the prices fixed; feedback
changes those prices and is treated below.

\section{How many sessions to admit}\label{L6:sec:admit}

The admission rule of \cref{L5:sec:scheduler} applies with two changes. First, each
instance has its own slot queue: with exponential holding times the wait for a slot is
the Erlang C wait of \cref{prop:mmm}, which grows without bound as the offered load of
slots approaches $m_P$ or $m_D$. Second, the Campbell estimate of resident KV
(\cref{cor:resident-kv}) is compared with the prefix cache of the prefill instance,
and memory is priced at $\theta_P$.

\section{Misses that feed themselves}\label{L6:sec:feedback}

So far the hit rate was an input. It is also an output. A miss re-inserts a whole
context into the prefill instance's pool, which evicts other paused prefixes sooner
(\emph{pool channel}); and a turn that waits longer is absent longer, which exposes its
own prefix to eviction (\emph{wait channel}). Let $W(h)$ be the mean wait at the prefill
instance at hit rate $h$, nonincreasing in $h$ (\cref{cor:mixture}), and let $G(T;h)$ be
the probability that a paused prefix survives an absence $T$ when the instance runs at
hit rate $h$, nonincreasing in $T$ and nondecreasing in $h$.

\begin{definition}[Feedback map]\label{L6:def:feedback}
The \term{feedback map} is $H(h)=\E_Z[\,G(Z+W(h);h)\,]$ on $[0,1]$, and an
\term{equilibrium} is a hit rate with $H(h)=h$.
\end{definition}

\paragraph{The wait channel is an admission rule.} Whether a waiting turn's prefix
can be evicted is decided by the engine, not by the workload. An engine that
takes the prefix out of the cache only when it schedules the turn (vLLM admits a
waiting request at the start of an iteration with the budget left, and until then
its cached blocks stay evictable) has the wait channel. An engine that pins the
prefix when the turn arrives has survival $G(Z;h)$ instead of $G(Z+W(h);h)$, a larger
value for every $h$, but it holds the pinned bytes and so lowers $G$ for the other
paused prefixes; the comparison of \cref{L6:prop:tarski}(ii) applies only when the
first effect dominates. Pinning thus helps the pinned turn and hurts the other paused
sessions. Which effect wins depends on the survival laws, so it must be computed for a
given workload and memory size.

$H$ is nondecreasing: if $h\le h'$ then $W(h)\ge W(h')$, so
$G(Z+W(h);h)\le G(Z+W(h');h)\le G(Z+W(h');h')$ for every $Z$, and taking expectations
preserves the order. No continuity is needed for what follows.

\begin{proposition}[Extremal equilibria]\label{L6:prop:tarski}
Let $H:[0,1]\to[0,1]$ be nondecreasing.
\begin{enumerate}[label=(\roman*)]
\item $\bar h=\sup\{x:x\le H(x)\}$ is an equilibrium, and every equilibrium is at most
  $\bar h$; $\underline h=\inf\{x:H(x)\le x\}$ is an equilibrium, and every
  equilibrium is at least $\underline h$.
\item If $\tilde H:[0,1]\to[0,1]$ is nondecreasing and $H\le\tilde H$ pointwise, then
  $\bar h_H\le\bar h_{\tilde H}$ and $\underline h_H\le\underline h_{\tilde H}$.
\end{enumerate}
\end{proposition}

\begin{proof}
(i) Let $A=\{x\in[0,1]:x\le H(x)\}$; $0\in A$, so $\bar h=\sup A$ exists. For $x\in A$,
$x\le H(x)\le H(\bar h)$ because $x\le\bar h$ and $H$ is nondecreasing; so $H(\bar h)$ is
an upper bound of $A$ and $\bar h\le H(\bar h)$. Then $H(\bar h)\le H(H(\bar h))$, so
$H(\bar h)\in A$ and $H(\bar h)\le\bar h$. Hence $H(\bar h)=\bar h$. Every equilibrium lies
in $A$, so it is at most $\bar h$. The argument for $\underline h$ is the mirror image
with $B=\{x:H(x)\le x\}\ni1$ and infima.
(ii) $\bar h_H=H(\bar h_H)\le\tilde H(\bar h_H)$, so $\bar h_H$ lies in the set $A$ of $\tilde H$, and
$\bar h_H\le\bar h_{\tilde H}$. Likewise $H(\underline h_{\tilde H})\le\tilde H(\underline h_{\tilde H})=\underline h_{\tilde H}$
puts $\underline h_{\tilde H}$ in the set $B$ of $H$, so $\underline h_H\le\underline h_{\tilde H}$.
\end{proof}

A larger prefix cache raises $G$ and so $H$, and by (ii) it raises both extremal
equilibria. Other changes need care in a closed loop: a shorter think time or a faster
decode also raises the turn rate, and the net effect on $H$ has no sign in general.

\begin{proposition}[Collapse]\label{L6:prop:collapse}
In the open prefill instance take the wait of \cref{thm:pk}, $W(h)=\infty$ where
$\lambda\E[S_h]\ge1$, and $G(\infty;h)=0$.
\begin{enumerate}[label=(\roman*)]
\item If $\lambda\,\E[S^{\mathrm{miss}}]\ge1$, then $H(0)=0$, and every $h$ at which the
  instance is overloaded is mapped to 0, which is mapped to itself: all-miss attracts
  the overloaded hit rates.
\item If moreover $H(h_1)\ge h_1$ for some $h_1>0$, the least equilibrium is 0 and the
  greatest is at least $h_1$: the instance has two distinct extremal equilibria.
\end{enumerate}
\end{proposition}

\begin{proof}
(i) At $h=0$ every turn misses, the load is $\lambda\E[S^{\mathrm{miss}}]\ge1$, so
$W(0)=\infty$ and $H(0)=G(\infty;0)=0$. At an overloaded $h$, $H(h)=0$ in the same way,
and $H(0)=0$. (ii) 0 is an equilibrium, so $\underline h\le0$; and $h_1\in A$, so
$\bar h\ge h_1$ by \cref{L6:prop:tarski}.
\end{proof}

In a closed system the population bounds the wait and removes this mechanism for $H(0)=0$,
but not the pool channel: if re-inserted contexts evict every paused prefix within a
think time, $H(0)=0$ still. Thrashing, a collapse of this kind cured by limiting the
load, is classical in paging systems; the admission cap of \cref{L6:sec:admit} is its
cure here.

\begin{example}[Two equilibria and a cap]\label{L6:ex:bistable}
Take deterministic $S^{\mathrm{hit}}=0.5$\,s, $S^{\mathrm{miss}}=5$\,s, a deterministic think time $Z=7$\,s and the
wait channel alone, $G(T)=e^{-T/60\,\mathrm{s}}$, with the Pollaczek--Khinchine wait. At
$\lambda=0.25$/s the all-miss load is $1.25$, $H=0$ on $[0,2/9]$, and a grid search finds
three equilibria: 0 (stable), about 0.250 (unstable, slope $H'\approx12$) and about
0.882 (stable, slope 0.07). At $\lambda=0.18$/s the all-miss load is $0.9$ and the only
equilibrium is about 0.885. Lowering the load removes the collapse without moving the
good equilibrium much.
\end{example}

\begin{proposition}[The forced-miss multiplier]\label{L6:prop:multiplier}
Let $H(h)=\nu+kh$ near an equilibrium, with $\nu\ge0$ and $0\le k<1$, so $h_0=\nu/(1-k)$, and
force a fraction $\delta\in[0,1]$ of turns to miss, so the realised hit rate is
$(1-\delta)H(h)$. The new equilibrium is $h_\delta=(1-\delta)\nu/(1-(1-\delta)k)$, and
\[
  h_0-h_\delta=\frac{\delta\,\nu}{(1-k)(1-(1-\delta)k)}\ \ge\ \delta\,h_0 .
\]
\end{proposition}

\begin{proof}See \cref{L6:ex:multiplier}.\end{proof}

To first order in $\delta$ a forced miss converts a would-be hit with probability $h_0$
and induces $h_0k/(1-k)$ further misses on other turns, each with its own price, so its
marginal price is $h_0[\Phi_{\mathrm{forced}}+\tfrac{k}{1-k}\bar\Phi_{\mathrm{induced}}]$.
This cache multiplier multiplies the queue factor $(1-\rho_P)/(1-\rho_P')$ of
\cref{thm:price}; it does not replace it. A multiplier common to all sessions scales
every $u^P_i$ by the same factor: it leaves the eviction order unchanged and moves only
the comparisons with $\theta_P$ (keep or drop, admit).

\section{What is exact and what is approximate}\label{L6:sec:honest}

The list of \cref{L5:sec:honest} applies. In addition,
\cref{L6:prop:link-price,L6:prop:shadow,L6:prop:tarski,L6:prop:collapse,L6:prop:multiplier}
are exact given the model. The new modelling choices are the link as fair bandwidth
sharing (a transfer engine that sends in arrival order would add a head-of-line term
to \cref{L6:prop:link-price}) and a single survival law with the mean wait in the
feedback map. The link constants $x_0$, $\kappa$ and $B$ must be measured, and the
survival law $G$ estimated from the workload.

\section{Summary}
\begin{itemize}
\item A transfer miss is priced exactly by insensitivity and grows linearly in the context; a prefill miss grows quadratically and has a head-of-line term.
\item Each pool has its own shadow price, and a decode cache is worth only the transfers it saves.
\item The hit rate is also an output: overload can add an all-miss equilibrium, and an admission cap removes it.
\end{itemize}

\section{Exercises}

\begin{exercise}\label{L6:ex:paths}
Derive \eqref{L6:eq:admitP}--\eqref{L6:eq:feedback} from \cref{L6:def:path}. Why can
turn $k$ be admitted only after turn $k-1$, and why does it reach the prefill server in
arrival order? Which steps take zero time?
\emph{Hint:} admission is in arrival order, and a FIFO stage serves only its first turn.
\end{exercise}

\begin{exercise}\label{L6:ex:nocache}
Disable prefix caching in the model of \cref{L6:def:path}. Show from
\cref{L6:def:path} that every turn then has $H_k=0$, and compute $\E[S]$ from
\eqref{L6:eq:ES}. Which of the conditions \eqref{L6:eq:stable} becomes harder to meet?
\end{exercise}

\begin{exercise}\label{L6:ex:compare}
A session with context $K$ is paused at both instances. Using \cref{L6:prop:shadow},
show that for $K$ large enough it is kept at the prefill instance whenever it is kept at
the decode instance, if $\theta_P=\theta_D$ and $b>0$.
\emph{Hint:} $\Phi_i\ge\Delta S_i\ge bK^2/2$ grows quadratically, $\Phi^X_i$ linearly in $K$.
\end{exercise}

\begin{exercise}\label{L6:ex:multiplier}
Prove \cref{L6:prop:multiplier}: $h_\delta$ solves $h=(1-\delta)(\nu+kh)$. Then take $k=0.5$ and $h_0=0.8$. Force $\delta=0.1$ of the turns
to miss. Compute $h_\delta$ and compare the loss $h_0-h_\delta$ with the direct loss
$\delta h_0$. How does the ratio behave as $k\uparrow1$?
\emph{Hint:} the ratio is $1/(1-(1-\delta)k)$.
\end{exercise}
